% !TeX root = main.tex
装 Visual Studio 之前，先搞清楚一件事：\textbf{这一代（18.x）和上一代（17.x）的差别不只是版本号}。MSVC 编译器的版本号从这一代起和 IDE 版本号"离婚"了，新建项目的默认 C++ 标准从 C++14 提到了 C++20，主机系统要求也从 Windows 10 提到了 Windows 11。这一章把这些变化讲清楚，让你在下载安装包之前就知道自己要装什么、会得到什么。

\section{版本谱系：从 17.x 到 18.x}
\subsection{版本号怎么读}
Visual Studio 有两套版本号，很多人第一次看会懵：

\begin{longtable}{L{2.4cm}L{5.0cm}L{2.8cm}L{4.0cm}}
\toprule
\textbf{名称} & \textbf{形式} & \textbf{例子} & \textbf{含义} \\
\midrule
\endfirsthead
\toprule
\textbf{名称} & \textbf{形式} & \textbf{例子} & \textbf{含义} \\
\midrule
\endhead
\bottomrule
\endlastfoot
\textbf{产品名} & 年份 & Visual Studio 2026 & 市场名称，按年发布 \\
\textbf{内部版本号} & \texttt{18.<次版本>.<维护号>} & \code{18.10.3} & 稳定通道的正式版本 \\
\textbf{Insiders 版本号} & \texttt{18.<次版本> Insiders <构建号>} & \code{18.10 Insiders 2} & 预览通道的版本 \\
\textbf{MSBuild 版本} & \code{18.x} & MSBuild 18.x & 与 IDE 主版本号绑定 \\
\end{longtable}

对照上一代：Visual Studio 2022 的内部版本号是 \code{17.x}，Visual Studio 2019 是 \code{16.x}。所以看到 \code{18.x}，就是 VS 2026。

\begin{tipbox}[提示]
在 IDE 里按 \ui{帮助 > 关于 Microsoft Visual Studio} 可以看到完整的版本号与构建号。写作本教程时，稳定通道最新为 \textbf{18.10.3}（2026 年 9 月 29 日发布，构建号 12224.181）。
\end{tipbox}

\subsection{关键时间线}
\begin{longtable}{L{3.4cm}L{11.6cm}}
\toprule
\textbf{时间} & \textbf{事件} \\
\midrule
\endfirsthead
\toprule
\textbf{时间} & \textbf{事件} \\
\midrule
\endhead
\bottomrule
\endlastfoot
2025 年 9 月 9 日 & \textbf{Visual Studio 2026 Insiders 通道上线}，取代 VS 2022 时代的 "Preview 通道" \\
2025 年 11 月 11 日 & \textbf{正式发布（GA）}：版本 18.0.0，构建号 11205.157，随附 MSVC 14.50 与 \code{v145} 工具集 \\
2026 年 5 月 & 版本 18.6，随附 MSVC 14.51 \\
2026 年 9 月 29 日 & 版本 18.10.3（写作时的最新稳定版） \\
2026 年 11 月（预计） & MSVC 14.52 进入稳定通道，预计提供 \code{/std:c++23} \\
\end{longtable}

\subsection{与 Visual Studio 2022 的关系}
\begin{itemize}
  \item \textbf{2022 仍在服务期内}：VS 2022 会一直提供安全更新到 2032 年 1 月，不必急着迁移。
  \item \textbf{VS 2026 的服务期}：到 2028 年 11 月。
  \item \textbf{安装目录独立}：VS 2026 与 VS 2022 分别装在 \codeP{C:\textbackslash{}Program Files\textbackslash{}Microsoft Visual Studio\textbackslash{}18} 和 \codeP{C:\textbackslash{}Program Files\textbackslash{}Microsoft Visual Studio\textbackslash{}2022} 下的版本子目录里，两者可以并存。
  \item \textbf{项目文件不用改}：用 VS 2022 建的 \code{.vcxproj} 可以直接用 VS 2026 打开，\textbf{不需要升级项目文件格式}。只有在你主动"重定目标"时，才会改写 Windows SDK 版本、平台工具集（\code{v143} -> \code{v145}）和 MSVC 工具集版本。
  \item \textbf{扩展兼容}：VS 2026 \textbf{开箱支持 Visual Studio 2022 的扩展}，这是这一代特意做的工作，减少了迁移时的扩展断供问题。
\end{itemize}

\begin{warnbox}[注意]
从 VS 2026 起，年度版本采用\textbf{就地更新}（in-place）而不是并排安装。也就是说将来从 2026 升到下一个年度版本时，是原地升级，不会像过去那样在"程序和功能"里躺着两个 Visual Studio。
\end{warnbox}

\section{选择版本（SKU）}
\subsection{四种产品}
\begin{longtable}{L{2.6cm}L{4.2cm}C{1.8cm}L{5.6cm}}
\toprule
\textbf{产品} & \textbf{面向对象} & \textbf{是否免费} & \textbf{说明} \\
\midrule
\endfirsthead
\toprule
\textbf{产品} & \textbf{面向对象} & \textbf{是否免费} & \textbf{说明} \\
\midrule
\endhead
\bottomrule
\endlastfoot
\textbf{Community} & 学生、个人开发者、开源项目、小型团队 & [OK] & 功能最全的免费版；\textbf{只能在 Stable 通道}，且必须在 180 天内更新到较新版本 \\
\textbf{Professional} & 专业开发者、中小团队 & [X] & 增加企业级协作与部分工具能力 \\
\textbf{Enterprise} & 大型企业 & [X] & 增加高级测试、诊断与架构工具（如代码覆盖率、IntelliTrace 等） \\
\textbf{Build Tools} & CI 服务器、无界面构建机 & [OK] & \textbf{只有编译器与构建工具，没有 IDE}；用于自动化构建流水线 \\
\end{longtable}

对绝大多数读者来说，\textbf{Community 就是正确选择}——C++ 开发所需的一切（MSVC 编译器、调试器、性能探查器、CMake 支持）它都有。

\begin{tipbox}[提示]
如果你只是想在服务器上跑 \code{msbuild} 编译，不要装完整 IDE，装 \textbf{Build Tools} 就够了，体积小得多。
\end{tipbox}

\subsection{通道：Stable 与 Insiders}
这一代把通道名称改了，很容易和"预览功能"混淆：

\begin{longtable}{L{2.4cm}L{3.2cm}L{3.4cm}L{5.2cm}}
\toprule
\textbf{通道} & \textbf{旧称} & \textbf{更新频率} & \textbf{适用场景} \\
\midrule
\endfirsthead
\toprule
\textbf{通道} & \textbf{旧称} & \textbf{更新频率} & \textbf{适用场景} \\
\midrule
\endhead
\bottomrule
\endlastfoot
\textbf{Stable} & Current（正式版） & 每月一次功能更新 & 日常开发、生产环境 \\
\textbf{Insiders} & Preview（预览版） & 更频繁，可提前约一个月拿到新工具集 RC & 尝鲜新特性、提前验证工具集 \\
\end{longtable}

需要注意：\textbf{"通道"和"预览功能"是两套独立机制}。即使你用的是 Stable 通道，也可以在 \ui{工具 > 选项 > 预览功能} 里逐个打开某些还没正式发布的单项功能。

\begin{warnbox}[注意]
Community 版\textbf{不能}加入 Insiders 通道。想尝鲜 Insiders，要么用 Professional/Enterprise，要么用 Insiders 的 Build Tools。
\end{warnbox}

\subsection{授权提醒}
Community 的免费授权有适用范围（个人、开源、学术，以及规模较小的组织）。团队使用前请核对官方授权条款——本教程不构成法律意见。

\section{MSVC 工具集：本代最大的架构变化}
\subsection{版本解耦}
这是理解 VS 2026 的关键。\textbf{从 VS 2026 起，MSVC 编译器的版本号不再与 Visual Studio 的版本号同步。}

过去的规律是：VS 2022 的每个小版本都对应一个 MSVC 版本（17.10 -> 14.40，17.11 -> 14.41……），一个萝卜一个坑。现在不一样了：MSVC 有了自己的节奏，\textbf{大约每 6 个月发布一个新版本}，而 IDE 仍然是每月更新一次。于是同一个 VS 版本里，你可以选择不同的工具集。

\begin{longtable}{L{2.2cm}C{1.8cm}C{2.0cm}L{3.8cm}L{4.0cm}}
\toprule
\textbf{MSVC 版本} & \textbf{\code{_MSC_VER}} & \textbf{编译器版本} & \textbf{首次随附于} & \textbf{支持类型与终止} \\
\midrule
\endfirsthead
\toprule
\textbf{MSVC 版本} & \textbf{\code{_MSC_VER}} & \textbf{编译器版本} & \textbf{首次随附于} & \textbf{支持类型与终止} \\
\midrule
\endhead
\bottomrule
\endlastfoot
\textbf{14.50} & \code{1950} & 19.50 & VS 2026 18.0（2025 年 11 月） & \textbf{长期支持（LTS）}，2028 年 11 月 \\
\textbf{14.51} & \code{1951} & 19.51 & VS 2026 18.6（2026 年 5 月） & 标准支持，2027 年 2 月 \\
\textbf{14.52} & \code{1952} & 19.52 & 预览通道 & 预览，正式发布后确定 \\
\end{longtable}

作为对比，VS 2022 最后几个版本是 14.43（\code{_MSC_VER} 1943）、14.44（\code{1944}）。\textbf{\code{_MSC_VER} 从 19xx 跳到 1950}，正是这次解耦的直接体现。

\begin{keybox}[实战场景]
团队里有人用 14.50、有人用 14.51，编出来的二进制能不能混用？可以——见下一节。
\end{keybox}

\subsection{平台工具集 \code{v145}}
在 MSBuild 项目（\code{.vcxproj}）里，编译器不是按版本号指定的，而是按\textbf{平台工具集}（Platform Toolset）指定的：

\begin{lstlisting}[style=textstyle]
VS 2015 -> v140
VS 2017 -> v141
VS 2019 -> v142
VS 2022 -> v143
VS 2026 -> v145   <- 本代
\end{lstlisting}

所以打开一个老项目，如果属性页里写着 \code{v143}，说明它还在用 VS 2022 的工具集。VS 2026 可以\textbf{向下兼容地}安装并使用旧工具集（最早可到 VS 2015），方便你逐个迁移项目。

工具集的实际安装位置：

\begin{lstlisting}[style=textstyle]
C:\Program Files\Microsoft Visual Studio\18\<版本>\VC\Tools\MSVC\
+-- 14.50.35717\     MSVC 14.50（长期支持）
+-- 14.51.36231\     MSVC 14.51

其中 14.50 / 14.51 是版本号，35717 / 36231 是具体工具集版本
\end{lstlisting}

\subsection{二进制兼容性}
\textbf{MSVC 14.50 与 14.51 都保持了与 VS 2015 及以后版本编译出的代码的二进制兼容性。}这条规则很重要：

\begin{itemize}
  \item 用 14.50 编译的\textbf{静态库}，可以链接进用 14.51 编译的程序；
  \item 用不同 MSVC 版本编译的 DLL 与 EXE 可以互相调用；
  \item 但\textbf{前提是运行库（CRT）链接方式一致}（都用 \code{/MD} 或都用 \code{/MT}）——这条规矩没变，详见\textbf{第 12 章}。
\end{itemize}

\begin{warnbox}[注意]
二进制兼容不代表行为一致。不同工具集版本在优化、警告、标准库实现细节上仍可能有差异，正式发布前应当用\textbf{固定的工具集版本}重新构建一次。
\end{warnbox}

\subsection{精确锁定工具集版本}
大型项目常常需要"锁死编译器版本"以保证可复现构建。三种做法：

\begin{lstlisting}[style=shellstyle]
:: 1) 命令行指定工具集版本
msbuild MyApp.sln /p:Configuration=Release /p:VCToolsVersion=14.50.35717

:: 2) 开发者命令提示符里选择工具集
vcvars64.bat -vcvars_ver=14.50

:: 3) 在项目或 Directory.Build.props 里固定
\end{lstlisting}

\begin{lstlisting}[style=xmlstyle]
<PropertyGroup>
  <VCToolsVersion>14.50.35717</VCToolsVersion>
</PropertyGroup>
\end{lstlisting}

如果这个被固定的工具集\textbf{装在另一个 Visual Studio 安装实例里}，从 VS 2026 18.8 起可以打开跨安装发现：

\begin{lstlisting}[style=xmlstyle]
<PropertyGroup>
  <EnableVCToolsVersionDiscovery>true</EnableVCToolsVersionDiscovery>
  <VCToolsVersion>14.50.35717</VCToolsVersion>
</PropertyGroup>
\end{lstlisting}

\section{C++ 语言与标准库现状}
\subsection{一个必须澄清的坑：默认标准有两个}
初学者最容易踩的坑在这里——\textbf{"默认 C++ 标准"有两个不同的答案，而且都对}：

\begin{longtable}{L{6.8cm}C{2.0cm}L{5.8cm}}
\toprule
\textbf{场景} & \textbf{默认标准} & \textbf{原因} \\
\midrule
\endfirsthead
\toprule
\textbf{场景} & \textbf{默认标准} & \textbf{原因} \\
\midrule
\endhead
\bottomrule
\endlastfoot
\textbf{新建}控制台应用 / Windows 桌面应用 / 动态库 / 静态库项目 & \textbf{C++20} & VS 2026 的项目\textbf{模板}在项目属性里写入了 \code{/std:c++20} \\
\textbf{编译器}在没有指定 \code{/std} 时的行为 & \textbf{C++14} & MSVC 编译器的历史默认值，为兼容旧代码而保留 \\
\end{longtable}

换句话说：新建项目你拿到的是 C++20；但如果你手写 \code{cl.exe} 命令、或者改了一个没设 \code{/std} 的老项目，你拿到的是 C++14。

\begin{tipbox}[提示]
不要依赖默认值。在 \ui{项目 > 属性 > C/C++ > 语言 > C++ 语言标准} 里显式选一个，团队协作时尤其如此。
\end{tipbox}

\subsection{可用的标准档位}
\begin{longtable}{L{4.8cm}L{3.8cm}L{6.0cm}}
\toprule
\textbf{开关} & \textbf{属性页选项} & \textbf{说明} \\
\midrule
\endfirsthead
\toprule
\textbf{开关} & \textbf{属性页选项} & \textbf{说明} \\
\midrule
\endhead
\bottomrule
\endlastfoot
\code{/std:c++14} & ISO C++14 标准 & 编译器的历史默认值 \\
\code{/std:c++17} & ISO C++17 标准 & \\
\code{/std:c++20} & ISO C++20 标准 & \textbf{新建项目的默认值}；自动启用 \code{/permissive-} \\
\code{/std:c++23preview} & 预览 - ISO C++23 标准 & 只启用 C++23 范围内的特性，不含实验性内容；\textbf{预览特性可能不保证 ABI 稳定} \\
\code{/std:c++latest} & 预览 - 最新 C++ 工作草案 & 启用所有已实现的新特性，含进行中的实验特性；自动启用 \code{/permissive-} \\
\end{longtable}

两个细节：

\begin{enumerate}
  \item \textbf{\code{/std:c++20} 与 \code{/std:c++latest} 会自动打开 \code{/permissive-}}（严格一致性模式），这会拒绝一些微软扩展语法，老代码可能因此报错。
  \item \textbf{检测当前标准用 \code{_MSVC_LANG}}，不要用 \code{__cplusplus}——除非你显式加了 \code{/Zc:__cplusplus}，否则 \code{__cplusplus} 永远是 \code{199711L}。
\end{enumerate}

\begin{lstlisting}[style=cppstyle]
#include <iostream>

int main() {
    std::cout << "_MSVC_LANG = " << _MSVC_LANG << '\n';
#if _MSVC_LANG >= 202302L
    std::cout << "当前至少是 C++23 档位\n";
#elif _MSVC_LANG >= 202002L
    std::cout << "当前是 C++20 档位\n";
#else
    std::cout << "当前是 C++17 或更低档位\n";
#endif
}
\end{lstlisting}

\subsection{C++23 与 C++26}
\textbf{C++23}：MSVC 14.50 起显著提高了 C++23 一致性，例如 \code{#warning} 预处理指令、\code{auto(x)} 退化拷贝等；14.51 继续补齐，并新增了标准库头文件 \textbf{\code{<flat_map>}} 与 \textbf{\code{<flat_set>}}。

\begin{warnbox}[注意]
写作本教程时（2026 年 10 月），稳定通道提供的 C++23 档位仍是 \textbf{\code{/std:c++23preview}}。据 C++ 团队博客，MSVC 14.52（预计 2026 年 11 月进入稳定通道）将提供正式开关 \textbf{\code{/std:c++23}}，届时 C++23 支持达到完整。属性页里的"预览 - ISO C++23 标准"选项会在新开关落地后自动指向 \code{/std:c++23}。
\end{warnbox}

\textbf{C++26}：截至写作时，官方文档\textbf{没有列出 C++26 核心语言特性的支持状态}，涉及 C++26 的公开内容主要是 constexpr \code{cmath}（\code{sin}/\code{cos}/\code{tan}/\code{pow}，通过 \code{/Zc:cmath} 启用）。想用前沿特性就挂 \code{/std:c++latest}，但要做好"下个版本可能变"的心理准备。

\section{面向 C++ 的 IDE 新特性速览}
以下都是 VS 2026 相对 VS 2022 的 C++ 相关变化，后续章节会逐个展开。

\subsection{界面与编辑}
\begin{itemize}
  \item \textbf{全新 UI，内置 11 套主题}，并可在 \ui{环境 > 视觉体验 > 主题颜色} 里直接改 Fluent 颜色令牌，不用装扩展。
  \item \textbf{开箱支持 VS 2022 扩展}。
  \item \textbf{C++ 属性（attributes）语法高亮}：\code{[[nodiscard]]}、\code{[[maybe_unused]]} 这类属性终于有自己的颜色了。
  \item \textbf{类视图支持子串搜索}，并大幅提升了大型项目的性能。
  \item \textbf{编辑器右下角增强}：可以直接查看和\textbf{修改文件编码}——这解决中文乱码问题时非常有用。
  \item \textbf{查找与替换支持排除文件}（Quick Find / Find in Files），大型仓库里搜索终于能排除 \code{third_party/} 了。
  \item \textbf{快捷键向其他编辑器看齐}：这一代统一了一批按键约定，从 VS Code 迁过来的人会舒服很多。
\end{itemize}

\subsection{项目系统与构建}
\begin{itemize}
  \item \textbf{右键直接生成预处理输出}：对着一个 \code{.cpp} 文件右键选 \textbf{预处理}，立刻得到预处理后的结果，排查宏和头文件包含问题时极其高效。
  \item \textbf{自带 CMake 4.1.1}，并且 CMake 提供了 \textbf{\code{Visual Studio 18 2026} 生成器}，还支持新的 \textbf{SLNX} 解决方案格式——现在可以从 CMake 侧构建 Visual Studio 的 C++ 项目。
  \item \textbf{设置助手（Setup Assistant）}：打开老项目时，帮你把缺失的工具集装上，或者把项目重定目标到新工具集。
  \item \textbf{后端优化}：相比 17.14，在 Unreal Engine City Sample 的 RenderThread 基准上有最高约 \textbf{6\%} 的运行性能提升。
\end{itemize}

\subsection{调试与诊断}
\begin{itemize}
  \item \textbf{调试器内联显示函数返回值}：单步跳过函数调用后，直接看到返回值，不用再设监视。
  \item \textbf{工具栏直接设置命令行参数}：\code{.vcxproj}、CMake、Unreal 项目都支持，不再局限于游戏开发工作负荷。
  \item \textbf{AddressSanitizer 支持 ARM64}（预览）。
  \item \textbf{clang-tidy 配置增强}：可分配更多处理器并行分析，也可追加自定义命令行参数（\ui{项目属性 > 代码分析 > Clang-Tidy}）。
  \item \textbf{容器调试支持 Podman}，与 Docker 体验一致。
  \item \textbf{采样式 PGO（SPGO）预览}、\textbf{Arm SVE} 支持、\textbf{Intel APX} 预览、PDB 体积裁剪。
\end{itemize}

\subsection{AI 辅助}
\begin{itemize}
  \item \textbf{Copilot Chat 的 Ask / Agent 两种模式}：Ask 回答问题，Agent 直接改代码。
  \item \textbf{编辑器右键新增 Copilot Actions}：把选中的代码丢给 Copilot，让它解释、优化、加注释、生成单元测试。
  \item \textbf{C++ 补全质量提升}：Copilot 会把相关文件作为上下文，减少胡编。
  \item \textbf{全代码库语义索引}：导航与着色性能提升 2 倍以上，并为 Copilot 提供跨文件的符号检索能力。
  \item \textbf{C++ 现代化（Modernization）正式可用}：用 \code{@Modernize} 之类的指令，让 Copilot 把 MSBuild 或 CMake 项目升级到更新的 MSVC，并顺手修掉警告与链接错误。
  \item \textbf{自带密钥（BYOK）}：可以接入 Microsoft Foundry、OpenAI、Anthropic、Ollama 等模型，不必登录 GitHub。
\end{itemize}

\begin{tipbox}[提示]
AI 功能是可选的。如果你所在的环境不允许联网或不允许上传代码，在 \ui{工具 > 选项 > GitHub > Copilot} 里关掉即可，IDE 的其他功能不受影响。
\end{tipbox}

\section{破坏性变更与迁移}
从 VS 2022 升级过来，下面这些必须知道。

\subsection{已移除的功能}
\begin{longtable}{L{4.2cm}L{4.6cm}L{5.8cm}}
\toprule
\textbf{移除项} & \textbf{影响} & \textbf{替代方案} \\
\midrule
\endfirsthead
\toprule
\textbf{移除项} & \textbf{影响} & \textbf{替代方案} \\
\midrule
\endhead
\bottomrule
\endlastfoot
\textbf{C++AMP} & GPU 并行扩展库没了 & 改用 CUDA / SYCL / OpenMP \\
\textbf{ARM32 工具链} & 不能再编译 ARM32 目标 & 用旧版 MSVC；新项目请转向 \textbf{ARM64} 或 ARM64EC \\
\code{/DEBUG:FASTLINK} & 链接器开关失效 & 改用 \textbf{\code{/DEBUG:FULL}} \\
\end{longtable}

\subsection{已弃用的功能}
\begin{itemize}
  \item \textbf{\code{/await} 编译开关}：将未来版本移除。C++14/17 项目想用标准协程，请用 \code{/await:strict}。
  \item \textbf{实验性协程头文件} \code{<experimental/coroutine>}、\code{<experimental/generator>}、\code{<experimental/resumable>}：请迁移到 C++20 的 \code{<coroutine>}。
  \item \textbf{MSVC 14.51 移除的一批老古董}：TR1（\code{std::tr1}）、\code{<hash_map>}/\code{<hash_set>}、\code{<experimental/filesystem>}、\code{stdext::checked_array_iterator} 等。这些大多已经弃用五年以上，但\textbf{老代码里确实还藏着}，升级时会集中爆出来。
  \item \textbf{工作负荷调整}：原 "C++ Linux" 工作负荷改名为 \textbf{"使用 C++ 的 Linux、Mac 和嵌入式开发"}；\textbf{使用 C++ 的移动开发（iOS/Android）与嵌入式/IoT 工具已不再受支持}，未来版本会移除。
\end{itemize}

\subsection{目标平台下限提高}
\textbf{MSVC 14.50 起，编译产物最低只能运行在 Windows 10 / Windows Server 2016 上}，不再支持 Windows 7、Windows 8、Windows 8.1。

请注意把两件事分开：

\begin{lstlisting}[style=textstyle]
主机（你运行 IDE 的机器）
  -> 必须是 Windows 11 或 Server 2019/2022/2025

目标（你编译出的 .exe 要运行的机器）
  -> 最低 Windows 10 / Server 2016
\end{lstlisting}

如果你的软件必须支持 Windows 7，那就不能升级到 14.50，只能停留在旧工具集上（并自行承担旧工具集停止支持的风险）。

\subsection{值得留意的新警告与行为变化}
\begin{itemize}
  \item \textbf{\code{/Zc:enumTypes}}：底层类型推导相关的新警告，升级后可能刷出一批。
  \item \textbf{P2266 更简洁的隐式移动}：可能影响重载决议，代码里可用 \code{__cpp_implicit_move} 做条件编译。
  \item \textbf{\code{<regex>} 多行模式行为变化}：默认改为标准行为；如需旧行为，把 \code{_REGEX_LEGACY_}\allowbreak\code{MULTILINE_MODE} 设为 \code{1}。同时 \code{<regex>} 在这一代做了大幅性能重写，\textbf{速度提升明显}。
\end{itemize}

\subsection{升级的推荐路径}
\begin{lstlisting}[style=textstyle]
1. 保留 VS 2022（别卸），装上 VS 2026
2. 用 VS 2026 直接打开老解决方案 -> 让它按 v143 编译，先确认能编过
3. 跑一遍测试，确认行为一致
4. 用"设置助手"或手动重定目标到 v145
5. 处理新出现的警告（/Zc:enumTypes、P2266 相关）
6. 用固定的 VCToolsVersion 重新构建并回归测试
7. 全部通过后，再考虑卸载 VS 2022
\end{lstlisting}

\section{我该装哪个？决策清单}
\begin{longtable}{L{6.2cm}L{8.8cm}}
\toprule
\textbf{你的情况} & \textbf{建议} \\
\midrule
\endfirsthead
\toprule
\textbf{你的情况} & \textbf{建议} \\
\midrule
\endhead
\bottomrule
\endlastfoot
学生 / 个人学习 C++ & \textbf{Community}，Stable 通道 \\
公司里写 Windows C++ 项目 & \textbf{Professional} 或 \textbf{Enterprise}（看公司授权） \\
服务器上做 CI 自动构建 & \textbf{Build Tools}，不要装 IDE \\
想提前验证新工具集 & \textbf{Insiders} 通道（注意 Community 不支持） \\
还在维护 Windows 7 兼容的软件 & 暂不升级到 14.50+，或双工具集并存 \\
主要写跨平台小工具 & 直接用 \textbf{VS Code}，不必装 VS \\
\end{longtable}

\section{小结}
\begin{itemize}
  \item VS 2026 的内部版本号是 \textbf{18.x}，2025 年 11 月 11 日正式发布，\textbf{主机要求 Windows 11}。
  \item \textbf{MSVC 版本与 IDE 版本解耦}：14.50（LTS，\code{_MSC_VER} 1950）、14.51（1951）、14.52（1952，预览），平台工具集统一为 \textbf{\code{v145}}。
  \item \textbf{新建 C++ 项目默认 C++20}，而编译器在未指定 \code{/std} 时默认仍是 \textbf{C++14}——两者别混。
  \item C++23 目前用 \textbf{\code{/std:c++23preview}}；预计 14.52 起提供正式的 \code{/std:c++23}。C++26 核心特性尚无公开支持。
  \item \textbf{已移除}：C++AMP、ARM32 工具链、\code{/DEBUG:FASTLINK}；\textbf{目标系统下限}提高到 Windows 10 / Server 2016。
  \item 日常开发选 \textbf{Community + Stable} 即可；CI 选 \textbf{Build Tools}。
\end{itemize}

\section{练习题}
\begin{enumerate}
  \item 打开 \ui{帮助 > 关于 Microsoft Visual Studio}，记下你机器上的完整版本号与构建号，并判断它属于 Stable 还是 Insiders。
  \item 写一个程序打印 \code{_MSC_VER}、\code{_MSVC_LANG}、\code{_MSC_FULL_VER}，分别在 \code{/std:c++14}、\code{/std:c++17}、\code{/std:c++20} 下编译，观察三者如何变化。
  \item 把一个 VS 2022 时代的项目用 VS 2026 打开，先\textbf{不}重定目标，确认能用 \code{v143} 编译通过；然后重定目标到 \code{v145}，记录出现了哪些新警告。
  \item 用 \code{/std:c++23preview} 写一段使用 \code{<flat_map>} 的代码（需要 14.51 及以上工具集），体会 C++23 标准库的新容器。
  \item 用 \code{msbuild /p:VCToolsVersion=...} 固定工具集版本构建一次，再用不固定的方式构建一次，对比两次输出的编译器版本。
\end{enumerate}

\section{相关笔记}
\begin{itemize}
  \item \textbf{本教程}：\textbf{Visual Studio 2026 使用教程}
  \item \textbf{下一章}：\textbf{安装与环境配置}
  \item \textbf{C++ 教程}：\textbf{C++ 教程}
  \item \textbf{VS 快速入门}：\textbf{Visual Studio 入门}
  \item \textbf{CMake}：\textbf{CMake 使用指南}
\end{itemize}

\section{参考链接}
\begin{itemize}
  \item \href{https://learn.microsoft.com/zh-cn/visualstudio/releases/2026/release-notes}{Visual Studio 2026 发行说明}
  \item \href{https://learn.microsoft.com/zh-cn/cpp/overview/what-s-new-for-msvc}{MSVC Build Tools 新增功能}
  \item \href{https://learn.microsoft.com/zh-cn/cpp/overview/compiler-versions}{MSVC 编译器版本对照表}
  \item \href{https://learn.microsoft.com/zh-cn/cpp/build/reference/std-specify-language-standard-version}{\code{/std} 语言标准开关}
  \item \href{https://learn.microsoft.com/en-us/visualstudio/releases/2026/vs-system-requirements}{Visual Studio 2026 系统要求}
  \item \href{https://learn.microsoft.com/zh-cn/cpp/porting/binary-compat-2015-2017}{C++ 二进制兼容性}
\end{itemize}
